\section{Đặc tả chi tiết các ca sử dụng cốt lõi (Use Case Specifications)}
\label{sec:use_case_specifications}

Dưới đây là đặc tả chi tiết của hai ca sử dụng nghiệp vụ trọng tâm trong hệ thống:

\begin{table}[htbp]
\centering
\small
\begin{tabularx}{\textwidth}{L{3.2cm} Y}
\toprule
\textbf{Mục đặc tả} & \textbf{Nội dung chi tiết} \\
\midrule
\textbf{Mã ca sử dụng} & \textbf{UC01} \\
\textbf{Tên ca sử dụng} & Kê đơn điện tử và kiểm tra tương tác thuốc CDSS \\
\textbf{Tác nhân chính} & Bác sĩ điều trị \\
\textbf{Tiền điều kiện} & Bác sĩ đã đăng nhập hệ thống với quyền hợp lệ và đã chọn hồ sơ bệnh nhân. \\
\textbf{Hậu điều kiện} & Đơn thuốc hợp lệ được lưu vào CSDL; các cảnh báo tương tác được ghi nhận vào nhật ký kiểm toán. \\
\midrule
\textbf{Luồng sự kiện chính (Basic Flow)} & 
\begin{enumerate}[leftmargin=1.2em, topsep=0pt, partopsep=0pt]
    \item Bác sĩ mở giao diện kê đơn CPOE từ danh sách buồng bệnh.
    \item Bác sĩ tìm kiếm và chọn thuốc từ danh mục 54.885 số đăng ký DAV.
    \item Bác sĩ nhập liều dùng, đường dùng, số lần dùng và thời gian điều trị.
    \item Hệ thống tự động kích hoạt Cổng CDSS Sentinel: kiểm tra trần liều 24 giờ và đối chiếu ma trận tương tác với các thuốc đã có trong đơn.
    \item Không phát hiện vi phạm tương tác hoặc vượt liều. Hệ thống hiển thị trạng thái an toàn (màu xanh).
    \item Bác sĩ bấm nút "Lưu đơn thuốc". Hệ thống lưu dữ liệu và phát thông báo thành công.
\end{enumerate} \\
\midrule
\textbf{Các luồng rẽ nhánh / ngoại lệ (Alternative Flows)} & 
\begin{itemize}[leftmargin=1.2em, topsep=0pt, partopsep=0pt]
    \item \textbf{4a. Phát hiện tương tác chống chỉ định (QĐ 5948):} Cổng Sentinel lập tức khóa nút Lưu đơn, đổi giao diện sang màu đỏ cảnh báo, hiển thị cơ chế nguy hiểm và kích hoạt luồng UC02 (Gợi ý thay thế 1-Click).
    \item \textbf{4b. Phát hiện bệnh nhân suy thận (eGFR $< 30$):} Hệ thống hiển thị cảnh báo vàng, đề xuất mức giảm liều theo Dược thư Quốc gia.
\end{itemize} \\
\bottomrule
\end{tabularx}
\caption{Bảng đặc tả ca sử dụng UC01: Kê đơn điện tử và kiểm tra tương tác CDSS}
\end{table}

\begin{table}[htbp]
\centering
\small
\begin{tabularx}{\textwidth}{L{3.2cm} Y}
\toprule
\textbf{Mục đặc tả} & \textbf{Nội dung chi tiết} \\
\midrule
\textbf{Mã ca sử dụng} & \textbf{UC02} \\
\textbf{Tên ca sử dụng} & Gợi ý thay thế thuốc BHYT tương đương dược lý 1-Click \\
\textbf{Tác nhân chính} & Bác sĩ điều trị \\
\textbf{Tiền điều kiện} & Đơn thuốc đang bị khóa do phát hiện tương tác chống chỉ định (từ bước 4a của UC01). \\
\textbf{Hậu điều kiện} & Thuốc vi phạm được thay thế bằng thuốc an toàn; nút Lưu đơn được mở khóa. \\
\midrule
\textbf{Luồng sự kiện chính (Basic Flow)} & 
\begin{enumerate}[leftmargin=1.2em, topsep=0pt, partopsep=0pt]
    \item Hệ thống kích hoạt thuật toán tìm kiếm thuốc thay thế 4 chiều: cùng nhóm ATC, được phép kê tại tuyến bệnh viện theo Thông tư 20/2022/TT-BYT, ưu tiên BHYT 100\%, và không tương tác với các thuốc còn lại.
    \item Hệ thống hiển thị thẻ gợi ý thuốc thay thế kèm lý do lâm sàng trên giao diện (ví dụ: Thay thế Esomeprazol bằng Pantoprazol 40mg).
    \item Bác sĩ bấm nút "Chấp nhận thay thế 1-Click".
    \item Hệ thống tự động gỡ bỏ thuốc vi phạm, thêm thuốc mới vào đơn, tính toán lại tương tác toàn đơn.
    \item Đơn thuốc đạt trạng thái an toàn, nút "Lưu đơn thuốc" được mở khóa trở lại.
\end{enumerate} \\
\bottomrule
\end{tabularx}
\caption{Bảng đặc tả ca sử dụng UC02: Gợi ý thay thế thuốc BHYT tương đương dược lý 1-Click}
\end{table}

\begin{table}[htbp]
\centering
\small
\begin{tabularx}{\textwidth}{L{3.2cm} Y}
\toprule
\textbf{Mục đặc tả} & \textbf{Nội dung chi tiết} \\
\midrule
\textbf{Mã ca sử dụng} & \textbf{UC03} \\
\textbf{Tên ca sử dụng} & Tra cứu y văn Dược thư RAG có đối soát số trang \\
\textbf{Tác nhân chính} & Bác sĩ điều trị, Dược sĩ lâm sàng \\
\textbf{Tiền điều kiện} & Người dùng đã đăng nhập hệ thống và mở giao diện buồng bệnh hoặc thư viện tài liệu. \\
\textbf{Hậu điều kiện} & Phản hồi được hiển thị theo mô hình hai phân vùng kèm con trỏ số trang có thể kiểm chứng. \\
\midrule
\textbf{Luồng sự kiện chính (Basic Flow)} & 
\begin{enumerate}[leftmargin=1.2em, topsep=0pt, partopsep=0pt]
    \item Người dùng nhập câu hỏi lâm sàng (ví dụ: liều dùng Ceftriaxon cho bệnh nhân suy gan, khoảng cách đưa thuốc).
    \item Bộ định tuyến NLU xác thực câu hỏi thuộc phạm vi dược lâm sàng hợp lệ.
    \item Module chuẩn hóa hoạt chất chuyển đổi danh từ thuốc sang tên INN chuẩn hóa.
    \item Động cơ truy xuất lai ghép (Hybrid Retrieval) tìm kiếm đoạn trích phù hợp nhất từ Dược thư Quốc gia.
    \item Mô hình ngôn ngữ lớn tổng hợp câu trả lời dựa trên ngữ cảnh được cung cấp, gán con trỏ trích dẫn số trang chính xác.
    \item Module hậu kiểm đối soát toàn bộ số liệu định lượng trong câu trả lời với văn bản gốc.
    \item Giao diện hiển thị tin nhắn phân tách hai vùng và cho phép bấm vào trích dẫn để mở trực tiếp trang PDF tương ứng.
\end{enumerate} \\
\midrule
\textbf{Các luồng rẽ nhánh / ngoại lệ (Alternative Flows)} & 
\begin{itemize}[leftmargin=1.2em, topsep=0pt, partopsep=0pt]
    \item \textbf{2a. Câu hỏi ngoài phạm vi y tế / cố ý can thiệp hệ thống:} Hệ thống từ chối trả lời và hướng dẫn quay lại ca bệnh.
    \item \textbf{4a. Điểm số tương đồng truy xuất thấp hơn ngưỡng tin cậy ($\theta < 0.72$):} Hệ thống thông báo không đủ căn cứ y văn và dẫn chiếu mở bản PDF đầy đủ.
    \item \textbf{6a. Phát hiện mâu thuẫn số liệu định lượng:} Kích hoạt cơ chế cắt mạch an toàn, thu hồi câu trả lời và cảnh báo kiểm tra tài liệu giấy.
\end{itemize} \\
\bottomrule
\end{tabularx}
\caption{Bảng đặc tả ca sử dụng UC03: Tra cứu y văn Dược thư RAG có đối soát số trang}
\end{table}

\begin{table}[htbp]
\centering
\small
\begin{tabularx}{\textwidth}{L{3.2cm} Y}
\toprule
\textbf{Mục đặc tả} & \textbf{Nội dung chi tiết} \\
\midrule
\textbf{Mã ca sử dụng} & \textbf{UC04} \\
\textbf{Tên ca sử dụng} & Quản lý nhân sự và xác thực phân quyền hệ thống \\
\textbf{Tác nhân chính} & Quản trị viên hệ thống \\
\textbf{Tiền điều kiện} & Quản trị viên đã đăng nhập thành công vào Cổng giám sát quản trị (\code{admin-portal}). \\
\textbf{Hậu điều kiện} & Hồ sơ nhân viên được lưu trữ trong CSDL SQLite; mật khẩu web được băm bảo mật theo chuẩn Argon2id. \\
\midrule
\textbf{Luồng sự kiện chính (Basic Flow)} & 
\begin{enumerate}[leftmargin=1.2em, topsep=0pt, partopsep=0pt]
    \item Quản trị viên mở phân hệ Quản lý nhân viên trên thanh điều hướng.
    \item Quản trị viên nhập thông tin nhân sự: Mã NV, Họ tên, Vai trò phân quyền, Khoa phòng, Số CCHN, Mã thẻ RFID, Mã PIN ca trực và Tên đăng nhập Web (@username).
    \item Quản trị viên tạo mật khẩu ban đầu hoặc sử dụng bộ sinh mật khẩu ngẫu nhiên an toàn.
    \item Hệ thống băm mật khẩu bằng thuật toán Argon2id ($m = 64\,\text{MB}, t = 3, p = 4$) và lưu bản ghi vào bảng \code{users}.
    \item Hệ thống ghi nhận sự kiện vào nhật ký kiểm toán mật mã (\code{AuditRecord}).
    \item Nhân viên có thể sử dụng tên đăng nhập và mật khẩu để truy cập Cổng lâm sàng hoặc Cổng quản trị theo đúng vai trò được cấp.
\end{enumerate} \\
\midrule
\textbf{Các luồng rẽ nhánh / ngoại lệ (Alternative Flows)} & 
\begin{itemize}[leftmargin=1.2em, topsep=0pt, partopsep=0pt]
    \item \textbf{2a. Trùng lặp tên đăng nhập Web (@username):} Hệ thống báo lỗi và gợi ý tên đăng nhập dựa trên họ tên không dấu.
    \item \textbf{4a. Nhân viên quên mật khẩu:} Quản trị viên thực hiện thao tác đặt lại mật khẩu; hệ thống xóa cache xác thực phiên của người dùng ngay lập tức và đánh dấu cờ \code{must_change_password = 1}.
\end{itemize} \\
\bottomrule
\end{tabularx}
\caption{Bảng đặc tả ca sử dụng UC04: Quản lý nhân sự và xác thực phân quyền hệ thống}
\end{table}
